\documentclass[preview]{standalone}

\usepackage{amsmath}
\usepackage{amssymb}
\usepackage{stellar}
\usepackage{definitions}

\begin{document}

\id{polynomials-over-rings}
\genpage

\section{Polynomials over rings}

\begin{snippet}{polynomial-rings-over-fields-summary}
    If \(K\) is a \field, then \(K[x]\) is a \euclideandomain and hence
    a \principalidealdomain and a \uniquefactorizationdomain. It is also
    a \noetherianring by the Hilbert basis theorem.

    More generally, for \(n\geq2\),
    \[
        K[x_1,\ldots,x_n]
        =
        K[x_1,\ldots,x_{n-1}][x_n].
    \]
    Repeated application of the Hilbert basis theorem shows that every
    \(K[x_1,\ldots,x_n]\) is Noetherian.
\end{snippet}

\begin{snippetproposition}{polynomial-ring-domain-and-units}{Domains and units in polynomial rings}
    Let \(A\) be an \integraldomain. For every \(n\geq1\), the \ring
    \[
        A[x_1,\ldots,x_n]
    \]
    is an \integraldomain and
    \[
        A[x_1,\ldots,x_n]^\times=A^\times.
    \]
\end{snippetproposition}

\begin{snippetproof}{polynomial-ring-domain-and-units-proof}{polynomial-ring-domain-and-units}{Domains and units in polynomial rings}
    We proceed by induction on \(n\). For \(n=1\), let
    \[
        f=a_0+a_1x+\cdots+a_rx^r,
        \qquad
        g=b_0+b_1x+\cdots+b_tx^t
    \]
    be nonzero elements of \(A[x]\), with \(a_r,b_t\neq0_A\). The
    leading coefficient of \(fg\) is \(a_rb_t\), which is nonzero because
    \(A\) is an \integraldomain. Thus
    \[
        \polynomialdeg(fg)=\polynomialdeg f+\polynomialdeg g
    \]
    and \(fg\neq0\), so \(A[x]\) is an \integraldomain.

    If \(fg=1_A\), then
    \[
        0=\polynomialdeg(fg)
        =\polynomialdeg f+\polynomialdeg g.
    \]
    Hence \(f=a_0\) and \(g=b_0\) are constant, with
    \(a_0b_0=1_A\). Consequently
    \[
        A[x]^\times=A^\times.
    \]

    Assume the result for \(n-1\). Since
    \[
        A[x_1,\ldots,x_n]
        =
        A[x_1,\ldots,x_{n-1}][x_n],
    \]
    the coefficient ring on the right is an \integraldomain and its
    units are precisely \(A^\times\). Applying the one-variable case
    gives
    \[
        A[x_1,\ldots,x_n]^\times
        =
        A[x_1,\ldots,x_{n-1}]^\times
        =
        A^\times,
    \]
    and proves that \(A[x_1,\ldots,x_n]\) is an \integraldomain.
\end{snippetproof}

\end{document}
